% Numerical-example section for the 3-D case (from the original ns3d project; see docs/navier_stokes/3d/NS3D_NOTES.md).
\documentclass[11pt]{article}
\usepackage[margin=1in]{geometry}
\usepackage{amsmath,amssymb}
\usepackage{booktabs}
\usepackage{siunitx}
\sisetup{round-mode=places,round-precision=2}

\newcommand{\Vh}{V_h}\newcommand{\Qh}{Q_h}
\newcommand{\Vdivh}{V_h^{\mathrm{div}}}
\newcommand{\eps}{\varepsilon}

\begin{document}

\section{Three-dimensional numerical verification}
\label{sec:3d-verification}

This section documents a three-dimensional numerical example for the Scott--Vogelius
BDF2--IMEX scheme of the preceding sections, run with the \texttt{ns3d} NumPy/SciPy
solver (Sections~\ref{sec:3d-mesh}--\ref{sec:3d-theorem2}) and, where noted, the DOLFINx
implementation \texttt{fenics\_bdf2\_imex\_sv\_3d.py}. The purpose of the example is to
confirm, on a smooth problem with a known exact solution, that the three-dimensional
implementation behaves as the theorems of the preceding sections predict: pressure-robust
velocity errors independent of the pressure (Section~\ref{sec:3d-pressure-robust}), the
expected spatial, temporal, and combined orders of convergence for the velocity (in the
$L^2$ and $H^1$ norms) and the pressure (in the $L^2$ norm)
(Sections~\ref{sec:3d-space}--\ref{sec:3d-combined}), and, directly, that the energy
stability inequality of Theorem~1 and the velocity error estimate of Theorem~2 hold
numerically with the constants and mechanisms their proofs claim
(Sections~\ref{sec:3d-theorem1}--\ref{sec:3d-theorem2}).

\subsection{Mesh}
\label{sec:3d-mesh}

The computational domain is the unit cube $\Omega=(0,1)^3$. For $n\in\mathbb N$, $\Omega$
is first partitioned into $n^3$ congruent sub-cubes; each sub-cube is split into $6$
tetrahedra by a Kuhn (corner-to-corner) triangulation, and each of the resulting
tetrahedra is further split into $4$ sub-tetrahedra by connecting its centroid to each of
its four faces (an Alfeld split). The mesh therefore consists of $24\,n^3$ tetrahedra, with
mesh size
\[
h=\frac{\sqrt3}{n}.
\]
The Alfeld refinement is not optional: for the Scott--Vogelius pair $\Vh/\Qh=P_k/P_{k-1}^{\mathrm{disc}}$
to be inf-sup stable in three dimensions one needs both the barycentric mesh structure and
polynomial degree $k\ge d=3$ (Zhang's theorem); the results below accordingly use the
minimal admissible pair $k=3$ (continuous, piecewise-cubic velocity) and $k-1=2$
(discontinuous, piecewise-quadratic pressure). Table~\ref{tab:3d-mesh} lists the mesh
sizes used for the spatial convergence study ($n=1,\dots,6$) and for the fixed mesh used
in the temporal convergence study ($n=10$); the pressure-robustness, combined-convergence,
and theorem-check studies of Sections~\ref{sec:3d-pressure-robust}
and~\ref{sec:3d-combined}--\ref{sec:3d-theorem2} use smaller, separately noted mesh lists
($n=1,2,3$) chosen for affordability rather than to match this table.

\begin{table}[h]\centering
\caption{Meshes used in Sections~\ref{sec:3d-space} and~\ref{sec:3d-time}: subdivisions
$n$ per side of the unit cube, mesh size $h=\sqrt3/n$, and number of tetrahedra $24n^3$.}
\label{tab:3d-mesh}
\begin{tabular}{crrr}
\toprule
study & $n$ & $h$ & \#tetrahedra \\
\midrule
space (Table~\ref{tab:3d-space}) & $1,2,3,4,5,6$ & $1.7321,\dots,0.2887$ & $24,\dots,5184$ \\
time (Table~\ref{tab:3d-time})   & $10$          & $0.3162$              & $24\,000$ \\
\bottomrule
\end{tabular}
\end{table}

\subsection{Manufactured solution}
\label{sec:3d-manufactured}

The scheme is verified against a smooth, exactly divergence-free velocity field with
homogeneous Dirichlet (no-slip) boundary data on $\partial\Omega$, together with a smooth
pressure field used to probe the pressure error. The velocity is constructed as the curl
of a vector potential, which is automatically divergence-free for any choice of potential:
let
\[
s(a)=a^2(1-a)^2,\qquad B(x,y,z)=s(x)\,s(y)\,s(z),\qquad
\mathbf A(x,y,z)=\begin{pmatrix}c_1\\ c_2\\ c_3\end{pmatrix}B(x,y,z),
\]
with $(c_1,c_2,c_3)=(1.0,\,0.7,\,-1.3)$, and set
\[
u(t,x,y,z)=g(t)\,\nabla\times\mathbf A(x,y,z),\qquad g(t)=1+\tfrac12\sin(2t).
\]
Since $s$ and $s'$ both vanish at $a=0,1$, every component of $\nabla\times\mathbf A$ --
each a linear combination of $\partial_xB,\partial_yB,\partial_zB$ -- vanishes on all six
faces of the cube, so $u=0$ on $\partial\Omega$ \emph{exactly}, for any choice of
$(c_1,c_2,c_3)$; the particular (asymmetric) values above were chosen only to avoid an
accidental extra symmetry in the test. The exact pressure is
\[
p(t,x,y,z)=R\cos(\pi x)\cos(\pi y)\cos(\pi z)\cos(t),
\]
where $R$ is a free amplitude used in Section~\ref{sec:3d-pressure-robust} to probe
pressure-robustness ($R=1$ in every other section). The right-hand side
$f=\partial_tu-\nu\Delta u+(u\cdot\nabla)u+\nabla p$ is obtained by exact symbolic
differentiation and enters the discrete scheme as an $L^2$-projected load.

Unless stated otherwise, every study in this section uses
\[
\nu=0.01,\qquad R=1,\qquad k=3\ (\text{i.e.\ }P_3/P_2^{\mathrm{disc}}),\qquad
\delta_1=10^{-3},\ \delta_2=\delta_3=0,
\]
i.e.\ only the first (and, in three dimensions, dominant) CIP jump term $S_1$ is active
(see the solver's scope note on why $S_2,S_3$ are not implemented in three dimensions);
volume integrals are evaluated with a quadrature rule exact enough for degree $qm=9$ per
coordinate direction, which was found necessary (at this value or higher) to keep
quadrature error from contaminating the pressure-robustness and convergence results on
this smooth-but-non-polynomial manufactured solution.

\subsection{Pressure-robustness}
\label{sec:3d-pressure-robust}

The exact pressure $p=R\cos(\pi x)\cos(\pi y)\cos(\pi z)\cos(t)$ is scaled by $R$, which
changes the forcing by $\nabla((R-1)p)$ but leaves the exact velocity completely
unchanged; a pressure-robust method's discrete velocity should not move either. Table
\ref{tab:3d-pressrob} reports, for $R\in\{1,10,100,1000,10000\}$ on a fixed mesh
($n=2$, $h\approx0.866$) with $\nu=0.01$, $\Delta t=0.05$, $T=0.2$: the $L^2$ norm of the
computed velocity, the maximum pointwise drift from the $R=1$ solution, and the maximum
pointwise divergence.

\begin{table}[h]\centering
\caption{Pressure-robustness: $k=3$, $n=2$ ($h\approx0.866$), $\nu=0.01$, $\Delta t=0.05$,
$T=0.2$. ``drift'' $=\max|u_h^R-u_h^{R=1}|$.}
\label{tab:3d-pressrob}
\begin{tabular}{rccc}
\toprule
$R$ & $\|u_h\|_{L^2}$ & drift & $\max|\nabla\cdot u_h|$ \\
\midrule
1     & $5.905801\times10^{-4}$ & $0$                      & $1.53\times10^{-12}$ \\
10    & $5.905801\times10^{-4}$ & $1.94\times10^{-12}$     & $1.36\times10^{-11}$ \\
100   & $5.905801\times10^{-4}$ & $2.11\times10^{-11}$     & $1.50\times10^{-10}$ \\
1000  & $5.905801\times10^{-4}$ & $2.13\times10^{-10}$     & $1.25\times10^{-9}$  \\
10000 & $5.905801\times10^{-4}$ & $2.14\times10^{-9}$      & $8.79\times10^{-9}$  \\
\bottomrule
\end{tabular}
\end{table}

The $L^2$ norm of the velocity is unchanged to the sixth significant digit across four
orders of magnitude in $R$, and the drift grows linearly with $R$ at a rate of
$\approx2\times10^{-13}$ per unit $R$ -- i.e.\ it is pure (quadrature and floating-point)
round-off, not a genuine sensitivity of the scheme to the pressure. This is the expected
pressure-robust behaviour of the Scott--Vogelius pair, and matches the two-dimensional
result (and the earlier finding, in that same two-dimensional work, that this clean a
result requires quadrature degree $qm\gtrsim9$ on this manufactured solution; $qm=6$ was
not sufficient there and was not tried again here).

\subsection{Spatial convergence}
\label{sec:3d-space}

Table~\ref{tab:3d-space} reports the velocity $L^2$ and $H^1$ errors and the pressure
$L^2$ error at $\bar t=T=0.2$, for $n=1,\dots,6$, with the time step fixed at the small
value $\Delta t=0.05$ so that the reported errors are dominated by the spatial
discretization rather than the time discretization. The ``order'' columns are the local
convergence rate
\[
\text{order}_i=\frac{\ln(e_{i-1}/e_i)}{\ln(h_{i-1}/h_i)}
\]
between consecutive rows (note $h_{i-1}/h_i=(i{+}1)/i$ here, \emph{not} a constant
factor of two, since $n$ is simply incremented by one each time -- a plain $\log_2$ ratio
would be wrong for every row but the first).

\begin{table}[h]\centering
\caption{Spatial convergence of the BDF2-IMEX Scott--Vogelius scheme ($P_3/P_2^{\mathrm{disc}}$),
$\nu=0.01$, $\Delta t=0.05$ (fixed), $T=0.2$.}
\label{tab:3d-space}
\begin{tabular}{rrcccccc}
\toprule
$n$ & $h$ & $\|u-u_h\|_{L^2}$ & order & $\|\nabla(u-u_h)\|_{L^2}$ & order & $\|p-p_h\|_{L^2}$ & order \\
\midrule
1 & 1.7321 & $4.6247\times10^{-4}$ & --   & $4.6074\times10^{-3}$ & --   & $6.5985\times10^{-2}$ & --   \\
2 & 0.8660 & $1.3919\times10^{-4}$ & 1.73 & $2.2223\times10^{-3}$ & 1.05 & $1.2313\times10^{-2}$ & 2.42 \\
3 & 0.5774 & $3.9949\times10^{-5}$ & 3.08 & $1.0013\times10^{-3}$ & 1.97 & $3.9107\times10^{-3}$ & 2.83 \\
4 & 0.4330 & $1.4344\times10^{-5}$ & 3.56 & $5.1034\times10^{-4}$ & 2.34 & $1.6909\times10^{-3}$ & 2.91 \\
5 & 0.3464 & $6.0806\times10^{-6}$ & 3.85 & $2.8798\times10^{-4}$ & 2.56 & $8.7567\times10^{-4}$ & 2.95 \\
6 & 0.2887 & $2.9532\times10^{-6}$ & 3.96 & $1.7632\times10^{-4}$ & 2.69 & $5.0992\times10^{-4}$ & 2.97 \\
\bottomrule
\end{tabular}
\end{table}

For the Scott--Vogelius pair $P_3/P_2^{\mathrm{disc}}$, the velocity error is expected,
asymptotically, to converge as $\|u-u_h\|_{L^2}=O(h^{k+1})=O(h^4)$ and
$\|\nabla(u-u_h)\|_{L^2}=O(h^{k})=O(h^3)$, and the pressure error as
$\|p-p_h\|_{L^2}=O(h^{k})=O(h^3)$. All three observed rates climb monotonically with
refinement and approach these targets from below: the $L^2$ velocity rate reaches $3.96$
by $n=6$ (target $4$), the pressure rate reaches $2.97$ (target $3$), and the $H^1$
velocity rate reaches $2.69$ (target $3$, the slowest of the three to converge, as is
typical since it measures one derivative more than the $L^2$ norm). This is the expected
pre-asymptotic behaviour of a cubic element on meshes this coarse ($n=1$ is only $24$
tetrahedra); no plateau or loss of order is visible over the range tested. Divergence
remained at $\mathtt{div\_linf}\le5.1\times10^{-9}$ at every mesh size, consistent with
the velocity staying (numerically) exactly divergence-free. Wall-clock cost grows sharply
with $n$ (from $1\,\si{s}$ at $n=1$ to $371\,\si{s}$ at $n=6$ for the full four-step run at
$\Delta t=0.05$), the expected consequence of $P_3$ elements on an unstructured tetrahedral
mesh with a sparse-direct solve repeated every time step.

\subsection{Temporal convergence}
\label{sec:3d-time}

Table~\ref{tab:3d-time} isolates the temporal error on a single, fixed, comparatively fine
mesh ($n=10$, $h\approx0.316$, $24\,000$ tetrahedra), by comparing each $\Delta t$'s
solution against a reference solution computed on the \emph{same} mesh at
$\Delta t_{\mathrm{ref}}=0.005$: since the spatial discretization is identical in every
comparison, the (fixed, and otherwise dominant on so fine a mesh) spatial error cancels
out of the difference, leaving the temporal error. The reference run used $N=40$ steps and
required approximately $11.6$ hours of wall-clock time; a fourth, finer point
($\Delta t=0.01$, $N=20$) was still in progress at the time of writing and is not included.

\begin{table}[h]\centering
\caption{Temporal convergence of the BDF2-IMEX scheme, $\nu=0.01$, fixed mesh $n=10$,
$T=0.2$, measured against a same-mesh reference solution at
$\Delta t_{\mathrm{ref}}=0.005$.}
\label{tab:3d-time}
\begin{tabular}{rrcccccc}
\toprule
$N$ & $\Delta t$ & $\|u-u_h\|_{L^2}$ & order & $\|\nabla(u-u_h)\|_{L^2}$ & order & $\|p-p_h\|_{L^2}$ & order \\
\midrule
2  & 0.08000 & $2.1368\times10^{-5}$ & --    & $1.7355\times10^{-4}$ & --    & $2.5317\times10^{-3}$ & --    \\
5  & 0.04000 & $2.1266\times10^{-7}$ & 6.65  & $1.7148\times10^{-6}$ & 6.66  & $4.1026\times10^{-9}$ & 19.24 \\
10 & 0.02000 & $5.0597\times10^{-8}$ & 2.07  & $4.0805\times10^{-7}$ & 2.07  & $9.8925\times10^{-10}$& 2.05  \\
\bottomrule
\end{tabular}
\end{table}

The BDF2--IMEX scheme is second order in time, so an asymptotic rate of $2$ is expected
in every column. This is exactly what is observed between $\Delta t=0.04$ and
$\Delta t=0.02$ ($2.05$--$2.07$ across all three quantities), but the first interval,
$\Delta t=0.08\to0.04$, shows an anomalously large apparent order ($6.65$, $6.66$, and
$19.24$ for the pressure). Taken at face value this is not a higher convergence order; it
indicates that the $\Delta t=0.08$ solution itself carries an unusually large error for
its step size, which is then resolved away by the time $\Delta t$ is halved to $0.04$,
rather than $\Delta t=0.04$ being unusually accurate. The most likely explanation is that
$\Delta t=0.08$ is close to, or beyond, a stability threshold associated with the
\emph{explicit} treatment of the convective term on this comparatively fine mesh
($n=10$): BDF2-IMEX schemes are only conditionally stable with respect to the explicit
part of the splitting, and the pressure (recovered through the discrete inf-sup
condition rather than solved for directly) is typically the most sensitive of the three
quantities to such a near-threshold effect, consistent with its far larger apparent order
over the same interval. This reading is a plausible interpretation of the data actually
obtained, not something separately confirmed by, e.g., a dedicated stability scan; the
clean, on-target rate recovered for every quantity between $\Delta t=0.04$ and $0.02$
indicates that the scheme is solidly inside its stable, asymptotic regime from
$\Delta t=0.04$ on. Including the pending $\Delta t=0.01$ point would help confirm
whether the rate continues to hold at $\approx2$ (as the trend between the last two points
suggests it should) or reveals a different pattern.

\subsection{Combined space--time convergence}
\label{sec:3d-combined}

Table~\ref{tab:3d-combined} refines $h$ and $\Delta t$ together, with
$\Delta t\approx0.15\,h^{7/4}$ chosen so that $\Delta t^4=O(h^7)$ stays subdominant to the
spatial term $\eps_h=(\nu+v_0h)h^{2k}=O(h^{2k+1})=O(h^7)$ that Theorem~2 predicts in the
convection-dominated regime $\nu\lesssim v_0h$ (here $\nu=0.01$ is indeed well below
$v_0h$ for every $h$ tested, with $v_0=1.5$) -- i.e.\ it is designed to let the claimed
\emph{superconvergent} rate $\|u-u_h\|_{L^2}=O(h^{k+1/2})=O(h^{3.5})$ show through, rather
than being masked by the temporal error. The resulting number of time steps grows quickly
with $n$ (3, 9, 17 for $n=1,2,3$), which is why only three mesh levels were affordable
here (the $n=3$ run alone took just over four minutes).

\begin{table}[h]\centering
\caption{Combined space--time convergence: $k=3$, $\nu=0.01$,
$\Delta t\approx0.15\,h^{7/4}$, $T=1$.}
\label{tab:3d-combined}
\begin{tabular}{rrrrcccccc}
\toprule
$n$ & $h$ & $N$ & $\Delta t$ & $\|u-u_h\|_{L^2}$ & order & $\|\nabla(u-u_h)\|_{L^2}$ & order & $\|p-p_h\|_{L^2}$ & order \\
\midrule
1 & 1.7321 & 3  & 0.3333 & $4.7946\times10^{-4}$ & --   & $9.8092\times10^{-3}$ & --   & $3.6367\times10^{-2}$ & --   \\
2 & 0.8660 & 9  & 0.1111 & $1.6942\times10^{-4}$ & 1.50 & $3.9712\times10^{-3}$ & 1.31 & $6.7883\times10^{-3}$ & 2.42 \\
3 & 0.5774 & 17 & 0.0588 & $4.9589\times10^{-5}$ & 3.03 & $1.6554\times10^{-3}$ & 2.16 & $2.1561\times10^{-3}$ & 2.83 \\
\bottomrule
\end{tabular}
\end{table}

The $L^2$ velocity rate climbs from $1.50$ to $3.03$ between the two available intervals,
consistent with pre-asymptotic approach to the superconvergent target $3.5$ -- not yet
there, but already past the non-superconvergent target $h^k=h^3$ that would be expected
without the pressure-robust cancellation, which is itself a meaningful confirmation of the
mechanism at only three mesh levels. The $H^1$ velocity and $L^2$ pressure rates ($2.16$
and $2.83$ at the finest interval) are likewise climbing towards their respective targets
of $h^k=h^3$. Divergence remained at $\mathtt{div\_linf}\le6.1\times10^{-12}$ throughout.
With only two intervals of data, climbing-but-not-yet-converged rates are the most this
experiment can honestly claim; confirming saturation at the asymptotic values would need
mesh levels beyond $n=3$, which were not affordable at this refinement rate in the time
available.

\subsection{Theorem 1 (energy stability) with forcing}
\label{sec:3d-theorem1}

Table~\ref{tab:3d-thm1} evaluates both sides of the energy stability inequality directly,
for several viscosities and one deliberately stressed configuration (forcing scaled by
$20$, initial data scaled by $3$), all at $n=2$, $\Delta t=0.1$, $T=0.5$. ``identity''
is the relative residual of the \emph{exact} identity
$E^{n+1}-E^n+\|u^{n+1}-2u^n+u^{n-1}\|^2+4\Delta t(\nu\|\nabla u^{n+1}\|^2+S^{n+1}+c)=
4\Delta t(f^{n+1},u^{n+1})$ (worst step); ``LHS/RHS'' is the ratio of the two sides of the
theorem itself, using the discrete dual norm $\|f\|_{-1,h}$ (a Poisson solve, as in the
two-dimensional verification); ``$|c|/\mathrm{diss.}$'' is the convective term relative to
the dissipation, which the theorem's proof requires to vanish exactly.

\begin{table}[h]\centering
\caption{Theorem 1 with forcing: $k=3$, $n=2$, $\Delta t=0.1$, $T=0.5$.}
\label{tab:3d-thm1}
\begin{tabular}{lrrrr}
\toprule
case & $\nu$ & identity & LHS/RHS & $|c|/\mathrm{diss.}$ \\
\midrule
$\nu=10^{-1}$                      & $10^{-1}$ & $2.85\times10^{-11}$ & $2.03\times10^{-5}$ & $3.55\times10^{-16}$ \\
$\nu=10^{-2}$                      & $10^{-2}$ & $3.10\times10^{-11}$ & $7.80\times10^{-7}$ & $8.80\times10^{-16}$ \\
$\nu=10^{-4}$                      & $10^{-4}$ & $6.92\times10^{-11}$ & $6.44\times10^{-9}$ & $5.05\times10^{-14}$ \\
stress: $f\times20$, $u_0\times3$  & $10^{-4}$ & $7.77\times10^{-10}$ & $1.23\times10^{-9}$ & $1.46\times10^{-12}$ \\
\bottomrule
\end{tabular}
\end{table}

The exact identity holds to $10^{-9}$--$10^{-11}$ in every case, including the stressed
configuration, and the convective term is below $5\times10^{-12}$ of the dissipation
throughout -- both confirm that the discrete velocity stays (numerically) exactly
divergence-free and that the pressure-cancellation / convective-cancellation mechanism
behind Theorem~1's proof holds in three dimensions exactly as it does in two. The
inequality itself holds with enormous room to spare in every case ($\mathrm{LHS/RHS}\ll1$),
and that room shrinks monotonically as $\nu$ decreases ($2\times10^{-5}$ at $\nu=10^{-1}$
down to $6\times10^{-9}$ at $\nu=10^{-4}$) -- exactly the $1/\nu$-type looseness already
identified and discussed for the two-dimensional version of this bound (the $H^{-1}$-type
pairing of the forcing term puts an explicit $1/\nu$ in the right-hand side), not a new
finding specific to three dimensions.

\subsection{Theorem 2: direct evaluation of the left-hand side}
\label{sec:3d-theorem2}

Table~\ref{tab:3d-thm2} evaluates the actual left-hand side of Theorem~2,
$\max_n\|e_h^n\|_0^2 + \nu\Delta t\sum\|\nabla e_h^{n+1}\|_0^2 + \Delta t\sum
S_{\mathrm{IMEX}}^{n+1}(e_h^{n+1},e_h^{n+1})$, using $e_h=u_h-s_h$ with $s_h$ the modified
Stokes projection of the exact solution (the same proxy for $w_h$ used in the
two-dimensional verification, since $w_h=\nabla\times P_h z$ itself is not available in
this code), on the same three runs as Table~\ref{tab:3d-combined} (so $n$, $h$, $N$,
$\Delta t$ are as there). ``$G_7$'' is $\max_n\|e_h^n\|_0/h^2$, the a~priori bound~(7)
check; ``$S(s_h)$'' is the stabilization-consistency term $\Delta t\sum
S_{\mathrm{IMEX}}^{n+1}(s_h^{n+1},s_h^{n+1})$; $\eps_h=(\nu+v_0h)h^{2k}$ with $v_0=1.5$.

\begin{table}[h]\centering
\caption{Left-hand side of Theorem 2 with $e_h=u_h-s_h$: $k=3$, $\nu=0.01$, same runs as
Table~\ref{tab:3d-combined}.}
\label{tab:3d-thm2}
\begin{tabular}{rrrrrrrr}
\toprule
$n$ & $G_7$ & $\|e_h\|^2$ & $\nu$-visc & $S(e_h)$ & $S(s_h)$ & LHS & LHS$/(\eps_h+\Delta t^4)$ \\
\midrule
1 & $3.85\times10^{-5}$ & $1.33\times10^{-8}$  & $1.28\times10^{-7}$ & $1.71\times10^{-12}$ & $2.32\times10^{-12}$ & $1.41\times10^{-7}$ & $2.00\times10^{-9}$ \\
2 & $3.94\times10^{-5}$ & $8.73\times10^{-10}$ & $6.43\times10^{-9}$ & $1.21\times10^{-13}$ & $4.51\times10^{-13}$ & $7.31\times10^{-9}$ & $1.32\times10^{-8}$ \\
3 & $2.96\times10^{-5}$ & $9.77\times10^{-11}$ & $1.25\times10^{-9}$ & $1.52\times10^{-14}$ & $8.68\times10^{-14}$ & $1.34\times10^{-9}$ & $4.14\times10^{-8}$ \\
\bottomrule
\end{tabular}
\end{table}

Three observations. First, the a~priori bound~(7) holds comfortably: $G_7=\max\|e_h\|/h^2$
stays around $3$--$4\times10^{-5}$ across all three meshes, far below any fixed constant
$\Gamma$ that would be needed for the bootstrap argument to close (Sec.~5.2 of the
corrected paper). Second, the left-hand side itself decreases rapidly with refinement
($1.4\times10^{-7}\to1.3\times10^{-9}$), dominated throughout by the viscous term
$\nu$-visc rather than the CIP dissipation $S(e_h)$, which is two to three orders of
magnitude smaller at every mesh size -- consistent with $S_1$ alone being a comparatively
mild stabilizer on this smooth solution, exactly as intended. Third, and worth stating
plainly rather than glossing over: the ratio of the left-hand side to the claimed scaling
$\eps_h+\Delta t^4$ is extremely small throughout ($\sim10^{-9}$ to $\sim10^{-8}$), several
orders of magnitude smaller than the corresponding ratios found in the two-dimensional
verification of this same bound (which ranged from a few tenths to several units). This
says the bound holds with very large room to spare here, not that it is wrong; part of
the explanation is structural -- $\eps_h\propto h^{2k}=h^6$ in three dimensions against
$h^{4}$ in two, and at the mesh sizes tested here $h$ ranges from $0.58$ to $1.73$, so this
extra power of $h$ changes the numerical scale of $\eps_h$ substantially relative to the
two-dimensional case -- but the ratio was not tracked finely enough across mesh levels in
this experiment to decompose that explanation quantitatively, and it should be treated as
an open observation rather than a settled one.

\end{document}
